We first establish consistency without imposing the smooth active-boundary
conditions needed for inference.

\begin{proposition}[Continuity and existence of the CML target]
\label{prop:population-existence}
Under Assumption~\ref{ass:basic-criterion}(1)--(3), $M$ is continuous on
$\Theta$ and attains its maximum.  Under part (4), the maximizer is the unique
target $\boldsymbol{\theta}^*$.
\end{proposition}

\begin{proof}
For $P_0$-almost every $\mathbf{x}$, the maximum of the finite collection of continuous
maps $\{q_j(\mathbf{x};\cdot):1\leq j\leq k\}$ is continuous.  The envelope in
Assumption~\ref{ass:basic-criterion}(3) and dominated convergence imply
continuity of $M$.  Existence follows from compactness of $\Theta$; uniqueness
is part (4).
\end{proof}

\begin{theorem}[Uniform convergence of the classification criterion]
\label{thm:uniform-convergence}
Under Assumption~\ref{ass:basic-criterion}(1)--(3),
\begin{equation}
 \sup_{\boldsymbol{\theta}\in\Theta}|M_n(\boldsymbol{\theta})-M(\boldsymbol{\theta})|\xrightarrow{p}0.
 \label{eq:uniform-convergence}
\end{equation}
If the observations form an infinite iid sequence, the convergence can be
taken almost surely.
\end{theorem}

\begin{proof}
The class $\{m_{\boldsymbol{\theta}}:\boldsymbol{\theta}\in\Theta\}$ has an integrable envelope.  It is
indexed by a compact finite-dimensional parameter set and is pointwise
continuous in $\boldsymbol{\theta}$ almost surely.  A finite covering argument, using
dominated continuity to control the oscillation within each cover element,
reduces the supremum to finitely many ordinary laws of large numbers.  This
gives \eqref{eq:uniform-convergence}; applying the strong law in the same
argument gives the almost-sure version.
\end{proof}

\begin{theorem}[Consistency of approximate CML maximizers]
\label{thm:cml-consistency}
Suppose Assumption~\ref{ass:basic-criterion} holds and
$\widehat{\boldsymbol{\theta}}_n$ satisfies \eqref{eq:approximate-maximizer}.  Then
\begin{equation}
 \widehat{\boldsymbol{\theta}}_n\xrightarrow{p}\boldsymbol{\theta}^*.
 \label{eq:cml-consistency}
\end{equation}
If the uniform convergence in Theorem~\ref{thm:uniform-convergence} and
$r_n\to0$ both hold almost surely, then the convergence in
\eqref{eq:cml-consistency} is almost sure.
\end{theorem}

\begin{proof}
Fix $\varepsilon>0$ and let
\[
 \Delta_\varepsilon=M(\boldsymbol{\theta}^*)-
 \sup_{\|\boldsymbol{\theta}-\boldsymbol{\theta}^*\|\geq\varepsilon}M(\boldsymbol{\theta})>0.
\]
On the event
\[
 2\sup_{\boldsymbol{\theta}\in\Theta}|M_n(\boldsymbol{\theta})-M(\boldsymbol{\theta})|+r_n
 <\Delta_\varepsilon,
\]
no approximate maximizer satisfying \eqref{eq:approximate-maximizer} can lie
outside the $\varepsilon$-ball around $\boldsymbol{\theta}^*$.  The probability of this
event tends to one by Theorem~\ref{thm:uniform-convergence} and $r_n=o_p(1)$.
The almost-sure statement follows identically.
\end{proof}

Theorem~\ref{thm:cml-consistency} deliberately separates statistical and
algorithmic claims.  A CEM iteration that terminates at an arbitrary local
maximum need not satisfy \eqref{eq:approximate-maximizer}.  Later we give a
local transfer result under consistent initialization and a quantified
optimization error.

For the first-order theory, Theorem~\ref{thm:population-curvature} supplies the
deterministic quadratic expansion
\begin{equation}
 M(\boldsymbol{\theta}^*+\mathbf{u})-M(\boldsymbol{\theta}^*)
 =\tfrac12\mathbf{u}^\mathsf THu+o(\|\mathbf{u}\|^2),
 \label{eq:population-quadratic-expansion}
\end{equation}
because $\nabla_{\boldsymbol{\theta}}M(\boldsymbol{\theta}^*)=0$ at an interior maximizer.  The next theoretical
layer will establish the matching local empirical-process expansion and show
that a sufficiently accurate local maximizer is asymptotically linear with
curvature $\mathbf{A}=\mathbf{I}_c-\mathbf{B}$.
